%% ma: L12-B08-0055
%% lop: 12
%% chuong: 2
%% bai: 8
%% dang: 12.8.30
%% loai: DS
%% muc: [NB, TH, VD, VD]
%% dap_an: "ĐĐSĐ"
%% nguon: "LaTeX giáo viên gửi 21/10/2026 (M05, Vectơ và hệ trục tọa độ trong không gian, 107 câu), Phần II Câu 10"
%% kien_thuc: "Bài 6 (quy tắc trung điểm, bình phương vô hướng), Bài 8 (trung điểm, khoảng cách hai điểm)"
%% trang_thai: chinh-thuc
%% ghi_chu: "Giáo viên đã duyệt 21/10/2026. Bỏ tiêu đề '(Bài toán Cực trị Tâm tỉ cự)'"
\begin{ex}
Trong không gian $Oxyz$, cho hai điểm cố định $A(2;0;-1)$, $B(4;6;1)$. Gọi $M(x;y;z)$ là một điểm tùy ý trong không gian.
\choiceTF
{\True Tọa độ trung điểm $I$ của đoạn thẳng $AB$ là $I(3;3;0)$.}
{\True $\overrightarrow{MA}+\overrightarrow{MB}=2\overrightarrow{MI}$.}
{Biểu thức $S=MA^2+MB^2$ đạt giá trị nhỏ nhất khi và chỉ khi $M$ trùng với điểm $A$.}
{\True Giá trị nhỏ nhất của biểu thức $S=MA^2+MB^2$ bằng $22$.}
\loigiai{a) $I(3;3;0)$. Đúng.

b) $\overrightarrow{MA}+\overrightarrow{MB}=2\overrightarrow{MI}+\overrightarrow{IA}+\overrightarrow{IB}=2\overrightarrow{MI}$. Đúng.

c) $S=(\overrightarrow{MI}+\overrightarrow{IA})^2+(\overrightarrow{MI}+\overrightarrow{IB})^2=2MI^2+IA^2+IB^2$ nhỏ nhất khi $M\equiv I$. Sai.

d) $AB^2=4+36+4=44$, $IA^2=IB^2=11$ nên $S_{\min}=22$. Đúng.}
\end{ex}
